\chapter{Configuration}
\label{app:config}

This appendix records the exact configuration used for the performance comparison in Section~\ref{sec:performance-comparison}. Both methods consume the same frozen, chronologically ordered data set and the same optimiser, so that any difference between them is attributable to the correlated Rao-Blackwellised construction rather than to data preparation.

\section{Data and split}

The data set is the frozen international football results file \texttt{results.csv}, containing $49{,}521$ rows. The rows are split chronologically into a training period $[1980\text{-}01\text{-}01, 2024\text{-}01\text{-}01)$, a test period $[2024\text{-}01\text{-}01, 2026\text{-}06\text{-}11)$, and a prediction period from 11 June 2026. The training set contains $4{,}702$ matches, the test set $305$, and the prediction set $104$. Matches with scores exceeding the truncation threshold $\max\_goals = 8$ are excluded before the split, and same-day repeats are retained as separate rows with zero elapsed time.

\section{Model parameters}

The static parameters are $\Theta = (\Gamma_0, B, \kappa, \alpha, \beta)$, with $\Gamma_0 \in \mathbb{R}^{K \times K}$, $B \in \mathbb{R}^{2 \times 2}$, $\kappa > 0$, $\alpha \in \mathbb{R}$ and $\beta \in \mathbb{R}$, subject to $\Gamma_0 \succ 0$, $B \succ 0$ and $\det(B) = 1$. The number of teams is $K = 48$. The initial mean is $\mu_0 = \mathbf{0}_{2K}$.

\section{Optimisation}

Both methods use Adam with a cosine-decay learning-rate schedule, learning rate $0.02$, seed $0$ and exactly $50$ updates. The SQMC log-likelihood gradient is averaged over $15$ independently randomised RQMC replicas at each update; the EKF gradient is deterministic. The EKF integrates the bivariate-Poisson likelihood under the predictive Gaussian by Gauss--Hermite quadrature of degree $32$.

\section{Full configuration}

Table~\ref{tab:config} lists the complete configuration. The values correspond to the authoritative comparison run \texttt{09092026\_1431} at source commit \texttt{336842f}, which is the run reported in Section~\ref{sec:performance-comparison}.

\begin{table}[H]
\centering
\begin{tabularx}{\textwidth}{lXX}
\toprule
Setting & Value & Description \\
\midrule
Number of teams $K$ & $48$ & Teams in the World Cup 2026 model \\
Particles $N$ & $512$ & SQMC particle count \\
Epochs & $50$ & Optimisation updates \\
Learning rate & $0.02$ & Adam step size \\
Learning-rate schedule & cosine decay & Decay of the learning rate over epochs \\
Seed & $0$ & Random seed for optimisation and RQMC \\
SQMC gradient replicas & $15$ & RQMC replicas averaged per gradient \\
Gauss--Hermite degree & $32$ & Quadrature degree for the EKF likelihood \\
$\max\_goals$ & $8$ & Score truncation threshold \\
$\text{match\_scale}$ & $1$ & Observation scale for non-friendly matches \\
Include friendlies & yes & Friendly fixtures retained in training \\
Training start & $1980\text{-}01\text{-}01$ & Start of the training period \\
Test start & $2024\text{-}01\text{-}01$ & Start of the test period \\
Prediction start & $2026\text{-}06\text{-}11$ & Start of the prediction period \\
GPU & A100 & Hardware for the SQMC run \\
Source commit & \texttt{336842f} & Repository revision of the run \\
Run ID & \texttt{09092026\_1431} & Identifier of the authoritative run \\
\bottomrule
\end{tabularx}
\caption{Configuration for the performance comparison of Section~\ref{sec:performance-comparison}.}
\label{tab:config}
\end{table}
